\documentclass[a4paper]{article}

% Español (México) con XeLaTeX: fontspec para fuentes Unicode, polyglossia
% para la división silábica y los nombres del documento en la variante mexicana.
\usepackage{fontspec}
\usepackage{polyglossia}
\setdefaultlanguage[variant=mexican]{spanish}
% Los nombres chinos van como «pinyin (original)»: xeCJK da a los caracteres
% Han su propia fuente; sin ella XeLaTeX los omite con un aviso.
% FandolSong no cubre algunos caracteres de nombres (U+73FA); WenQuanYi Zen Hei sí.
\usepackage[AutoFallBack=true]{xeCJK}
\setCJKmainfont{FandolSong-Regular.otf}
\setCJKfallbackfamilyfont{\CJKrmdefault}{WenQuanYi Zen Hei}
% polyglossia no traduce los nombres de listings.
\AtBeginDocument{\providecommand{\lstlistingname}{}\renewcommand{\lstlistingname}{Listado}%
  \providecommand{\lstlistlistingname}{}\renewcommand{\lstlistlistingname}{Índice de listados}}
\usepackage{amsmath,amssymb}
\usepackage{graphicx}
\usepackage[margin=2.5cm]{geometry}
\usepackage[most]{tcolorbox}
\usepackage{etoolbox}
\usepackage{subcaption}
\usepackage{listings}
\usepackage{hyperref}
\usepackage{booktabs}
\usepackage{float}

\newtcolorbox{knowledgebox}[1]{
    enhanced, colback=blue!5!white, colframe=blue!75!black, colbacktitle=blue!75!black,
    coltitle=white, fonttitle=\bfseries, title={#1}, boxrule=1pt, sharp corners
}

\newtcolorbox{importantbox}[1]{
    enhanced, colback=yellow!10!white, colframe=yellow!80!black, colbacktitle=yellow!80!black,
    coltitle=black, fonttitle=\bfseries, title={#1}, sharp corners
}

\newtcolorbox{warningbox}[1]{
    enhanced, colback=red!5!white, colframe=red!75!black, colbacktitle=red!75!black,
    coltitle=white, fonttitle=\bfseries, title={#1}, sharp corners
}

\lstset{
    basicstyle=\ttfamily\small,
    keywordstyle=\color{blue},
    stringstyle=\color{red!60!black},
    commentstyle=\color{green!60!black},
    breaklines=true,
    frame=single,
    numbers=left,
    numberstyle=\tiny\color{gray},
    captionpos=b,
    extendedchars=true
}

\newcommand{\notetitle}{CS224R Lecture 15: imitación y aprendizaje por refuerzo jerárquicos}
\newcommand{\noteauthors}{Elaborado a partir de la clase pública de Chelsea Finn en Stanford CS224R Lecture 15}
\newcommand{\notedate}{18 de marzo de 2025}
\newcommand{\videochannel}{Stanford Online}
\newcommand{\videopublishdate}{5 de marzo de 2025}
\newcommand{\videoduration}{aprox. 80 minutos}
\newcommand{\videourl}{https://cs224r.stanford.edu/}
\newcommand{\videocoverpath}{cover.jpg}
\newcommand{\repourl}{https://github.com/hqhq1025/ai-course-notes}


% Footer with repo info
\usepackage{fancyhdr}
\pagestyle{fancy}
\fancyhf{}
\fancyfoot[C]{\thepage}
\fancyfoot[R]{\footnotesize\color{gray}\href{\repourl}{AI Course Notes}}
\renewcommand{\headrulewidth}{0pt}
\renewcommand{\footrulewidth}{0pt}

\begin{document}
\begin{titlepage}
\centering
{\Large Notas del curso\par}
\vspace{1.2cm}
{\huge\bfseries \notetitle\par}
\vspace{0.8cm}
{\large \noteauthors\par}
\vspace{0.3cm}
{\large \notedate\par}
\vspace{1.1cm}

\ifdefempty{\videocoverpath}{
{\small No se proporcionó imagen de portada.\par}
}{
\includegraphics[width=0.82\textwidth,height=0.45\textheight,keepaspectratio]{\videocoverpath}\par
}

\vfill
\begin{tcolorbox}[width=0.9\textwidth, colback=black!2!white, colframe=black!60, sharp corners]
\textbf{Autor o canal del video}: \videochannel\par
\textbf{Fecha de publicación}: \videopublishdate\par
\textbf{Duración del video}: \videoduration\par
\textbf{Ponente}: Chelsea Finn\par
\textbf{Página del curso}: \href{https://cs224r.stanford.edu/}{cs224r.stanford.edu}\par
\textbf{Página del proyecto}: \href{\repourl}{\nolinkurl{\repourl}}\par
\end{tcolorbox}
\end{titlepage}

\tableofcontents
\newpage

%% --- Inicio del contenido --- %%

\section{Introducción y panorama del contenido}

\subsection{Ubicación de la clase y objetivos de aprendizaje}

En la Lecture 15, Chelsea Finn analiza a fondo el marco de ingeniería y las vías de gobernanza de la imitación y el aprendizaje por refuerzo jerárquicos. Esta sección no solo busca que entendamos por qué la estructura jerárquica es indispensable en las tareas de horizonte largo. También busca reproducir el ciclo cerrado «exploración estructurada → observación → gobernanza», que el ponente subraya una y otra vez a lo largo de la clase. La clase abarca datos, arquitectura, diseño de prompts y objetivos, y casos prácticos. Todas las notas se reorganizan según la lógica didáctica para facilitar su reproducción posterior.

\begin{knowledgebox}{Principios para construir notas con una jerarquía de dos niveles}
La clase se divide en: 1) motivación, 2) marco, 3) subobjetivos/imitación, 4) lenguaje + llamadas a LLM, 5) gobernanza de ingeniería. Cada bloque debe resaltar sus ideas centrales con `importantbox/knowledgebox/warningbox`. Las comparaciones se organizan con tablas o imágenes. Al final, un resumen sintetiza el contenido para facilitar el repaso.
\end{knowledgebox}

\subsection{Línea de tiempo y mapa de materiales}

\begin{table}[H]
\centering
\begin{tabular}{p{3.5cm}p{8cm}}
\toprule
Intervalo de tiempo & Contenido destacado y materiales \\
\midrule
00:00--00:12 & El ponente expone las razones de la jerarquía, long horizon y hierarchical induction bias (slide-01～05) \\
00:12--00:30 & Formalización de Options/SMDP y value estimation (slide-06～12) \\
00:30--00:45 & Jerarquía goal-conditioned; comparación entre HIRO y HAC (slide-13～20) \\
00:45--00:55 & Imitación jerárquica y descubrimiento de habilidades (slide-21～28) \\
00:55--01:12 & LLM + planificación con lenguaje, casos avanzados y recomendaciones de despliegue (slide-29～40) \\
01:12--01:20 & Checklists prácticas y cierre con Q\&A (slide-41～46) \\
\bottomrule
\end{tabular}
\caption{Línea de tiempo de la Lecture 15 y su correspondencia con las slides.}
\end{table}

\begin{figure}[H]
\centering
\includegraphics[width=0.85\textwidth]{slides-images/slide-01.jpg}
\caption{Slide de apertura de la clase: por qué merece atención la estructura jerárquica (slide-01).}
\end{figure}

\begin{importantbox}{Estrategia de uso de materiales}
Se citan primero las capturas de las slides. Si el contenido de una slide no basta, se complementa con los subtítulos o las notas. Cada imagen debe indicar su ruta dentro del directorio (por ejemplo, `slides-images/slide-xx.jpg`), y su caption debe incluir el timecode para poder ubicarla después.
\end{importantbox}

\subsection{Resumen de la sección}

Esta sección traza un mapa para repasar la clase a partir de su ubicación, la línea de tiempo y la guía de materiales. Todo el contenido posterior se desarrolla en este orden: motivación → modelo → práctica.
\section{Motivación jerárquica y descomposición de tareas}

\subsection{Trampas típicas del horizonte largo}

Ante recompensas dispersas, explosión combinatoria y entornos no estacionarios, una política de un solo nivel suele caer en el sesgo previo de «ajuste fino-sobreajuste»; la postura de Jason Punchline es no abandonar la perspectiva estructurada. Al agregar jerarquía, redefinimos la granularidad del espacio de acciones: cada action macro se encarga de una subtarea de varios pasos, lo que reduce la presión sobre la asignación de crédito.

\begin{knowledgebox}{Tres pasos para la descomposición estructurada en horizontes largos}
\begin{enumerate}
    \item Identificar los check-point clave (ejemplo: «sacar la olla» en una tarea de cocina);
    \item Extraer el subgoal que requiere cada check-point;
    \item Diseñar una subpolítica que cumpla ese subgoal y, al completarlo, devuelva el reward global.
\end{enumerate}
\end{knowledgebox}

\begin{warningbox}{El costo de ignorar la descomposición}
\begin{itemize}
    \item En una política plana, cualquier falla a lo largo de una trayectoria larga obliga a volver a explorar;
    \item El modelo de credit assignment se deja activar fácilmente por subtareas iniciadas por ruido, lo que provoca oscilaciones en el entrenamiento;
    \item Un agente sin jerarquía difícilmente puede comprimir el contexto y colapsa en entornos parcialmente observables.
\end{itemize}
\end{warningbox}

\subsection{Jerarquía y sesgo inductivo}

En cierta medida, la jerarquía también es un sesgo inductivo: le decimos al modelo «primero cumple una serie de objetivos intermedios con sentido y luego pasa al cierre». Inyectar correctamente este sesgo puede reducir de manera notable la sample complexity; pero si el sesgo es incorrecto (por ejemplo, un número fijo de skill o la falta de una stop condition general), se forma un training bottleneck.

\begin{table}[H]
\centering
\begin{tabular}{p{4cm}p{5cm}p{3cm}}
\toprule
Dimensión & RL tradicional & RL jerárquico \\
\midrule
Granularidad de decisión & Acciones atómicas & Subpolíticas/macroacciones \\
Eficiencia de muestreo & Baja, requiere rollout completo & Alta, las subpolíticas son reutilizables \\
Estrategia de exploración & Exploración pura & Construir subgoal reduce la search \\
Asignación de crédito & Retorno global & Reward adicional por subpolítica \\
\bottomrule
\end{tabular}
\caption{Comparación de la descomposición de tareas entre el RL tradicional y el RL jerárquico.}
\end{table}

\subsection{Plantillas jerárquicas en escenarios industriales}

En productos reales (como robots de cocina o plataformas de pruebas automatizadas), una plantilla jerárquica suele incluir: 1) módulo de percepción/análisis (extrae subgoal), 2) Strategy Planner (elige option/high-level plan), 3) Execution Layer (skill + verifiers), 4) Recovery + Logging (incident grade-book + prompt log). Cada capa debe registrar en texto claro su input/output en `notes/ops/`, para que el equipo on-call pueda cotejarlo con rapidez.

\begin{knowledgebox}{Plantillas industriales comunes}
\begin{itemize}
    \item Perception: state filtering + object detection;
    \item Planner: high-level policy combinada con options;
    \item Executor: skill library + environment interaction;
    \item Recovery: failure flag + fallback skill.
\end{itemize}
\end{knowledgebox}

\subsection{Resumen de la sección}

La parte de motivación subraya que lo «estructurado» no es un add-on, sino el núcleo para resolver la asignación de crédito y la exploration; las partes siguientes sobre Options, goal-conditioned e imitación se desarrollan en torno a estas dos líneas principales.
\section{Marco de options e implementación de SMDP}

\subsection{Formulación Semi-MDP y criterios de modelado}

El marco de options extiende el MDP a un Semi-MDP, lo que permite que la policy elija macro-actions que duran varios pasos. Denotamos con $\tau_\omega$ la longitud efectiva de ejecución de una option; el reward esperado se modela como $\mathbb{E}\left[\sum_{t=0}^{\tau_\omega} \gamma^t r_t\right]$, y el discounted factor conserva la convergencia.

\begin{importantbox}{La terna del diseño de una option}
\begin{itemize}
    \item $\mathcal{I}_\omega$: en qué estados puede activarse la option;
    \item $\pi_\omega$: la política interna se encarga de las acciones de varios pasos;
    \item $\beta_\omega$: la función de terminación controla el horizon de la option.
\end{itemize}
\end{importantbox}

\subsection{Estimación de valor y ecuación de Bellman}

La función value de una option se expresa como \(\mathrm{Q}(s, \omega) = \mathbb{E}\left[\sum_{t=0}^{\tau_\omega} \gamma^t r_t + \gamma^{\tau_\omega} V(s')\right]\). Esto nos permite aplicar Q-learning (o Actor-Critic) en la política de alto nivel sin tener que hacer actualizaciones combinatorias para cada acción atómica.

\begin{knowledgebox}{Desarrollo de Bellman semimarkoviano}
\(\quad V(s)=\sum_\omega \pi_{\text{high}}(\omega|s)\left[\mathbb{E}_{\tau_\omega}\left[\sum_{t=0}^{\tau_\omega}\gamma^t r_t\right]+\mathbb{E}[\gamma^{\tau_\omega}V(s')]\right]\)
\end{knowledgebox}

\begin{warningbox}{Sobreajuste de las options}
Si $\mathcal{I}_\omega$ es demasiado estrecho o la estructura de $\pi_\omega$ es demasiado rígida, la política de alto nivel cae en local minima del tipo «solo unas cuantas options son eficaces». Se recomienda aplicar regularization (dropout style) o ampliar dinámicamente las options disponibles.
\end{warningbox}

\begin{table}[H]
\centering
\begin{tabular}{p{4cm}p{5cm}}
\toprule
Componente & Función \\
\midrule
Initiation set & Controla cuándo puede el nivel alto elegir una option \\
Intra-option policy & Se encarga de la ejecución de varios pasos, junto con el skill replay \\
Termination function & Decide si la option se mantiene \\
High-level policy & Elige options con una granularidad temporal gruesa \\
\bottomrule
\end{tabular}
\caption{Componentes clave del marco de options.}
\end{table}

\subsection{Mantenimiento y ampliación del conjunto de options}

En proyectos grandes hay que mantener un conjunto de options: el skill engineer actualiza el initiation set y la función de termination según la coverage/usage matrix, y periódicamente retira las options con baja tasa de uso o con un failure rate alto. El usage log puede escribirse en un instrumentation dashboard para facilitar la governance review.

\begin{figure}[H]
\centering
\includegraphics[width=0.85\textwidth]{slides-images/slide-08.jpg}
\caption{Flujo de entrenamiento de options y usage tracking (slide-08).}
\end{figure}

\begin{importantbox}{Recomendaciones para operar el conjunto de options}
\begin{itemize}
    \item Usar una coverage matrix para supervisar cuántas veces se invoca cada option;
    \item Cada cambio en una option debe actualizar a la vez el prompt log y la instrumentation;
    \item Las options nuevas se prueban primero en un sandbox y después se pasan a production;
\end{itemize}
\end{importantbox}

\subsection{Resumen de la sección}

Options/SMDP aporta la formalización matemática del RL jerárquico: pone el acento en la extensión temporal de las action y en la reformulación de la función value, y en los proyectos reales exige registrar initiation/termination/logs en la instrumentation.
\section{Políticas jerárquicas Goal-Conditioned}

\subsection{Subobjetivos de alto nivel y ejecución de bajo nivel}

La política de alto nivel produce subobjetivos continuos o discretos \(g\), y la política de bajo nivel toma \((s, g)\) como conditioning para muestrear acciones. En el continuous case, el subobjetivo puede definirse como un embedding de alta dimensión de los elementos; en el discrete case, puede usarse directamente como el id de una habilidad.

\begin{table}[H]
\centering
\begin{tabular}{p{3cm}p{5cm}p{4cm}}
\toprule
Método & Mecanismo central & Tratamiento típico \\
\midrule
HIRO & Off-policy correction & El nivel alto restablece el objetivo cada $c$ pasos \\
HAC & Recursive hindsight relabeling & La reward del nivel alto proviene directamente de los subobjetivos \\
FuN & Predictive subgoal representations & Usa un predictive model para medir si el subgoal es alcanzable \\
\bottomrule
\end{tabular}
\caption{Comparación de métodos jerárquicos goal-conditioned comunes.}
\end{table}

\begin{knowledgebox}{Consideraciones para diseñar subobjetivos}
\begin{itemize}
    \item Alcanzabilidad: el subobjetivo debe poder lograrse dentro del horizon del nivel bajo;
    \item Identificabilidad: la política de bajo nivel necesita saber cuándo ha completado el subobjetivo;
    \item Reutilización: conviene diseñar subobjetivos generales para evitar que el nivel alto cree objetivos nuevos con frecuencia;
    \item Estabilidad: la distribución de subobjetivos de alto nivel debe ser acotada y se recurre al hindsight relabel para controlar el drift.
\end{itemize}
\end{knowledgebox}

\subsection{Retos del entrenamiento y ciclos de retroalimentación}

El entrenamiento del RL jerárquico suele venir acompañado de non-stationarity: la política de alto nivel se enfrenta a un comportamiento de bajo nivel que cambia constantemente, y la política de bajo nivel depende, a su vez, de subobjetivos dinámicos. El ponente recomienda usar un replay buffer de doble vía (uno para high y otro para low) y escribir un log que registre el hash del high-level goal y la associated reward.

\begin{warningbox}{Propagación del gradiente y conflictos en la repetición de experiencias}
\begin{itemize}
    \item el low-level replay puede sobreajustarse a los high-level plans tempranos;
    \item la high-level reward requiere un hindsight relabel basado en la low-level completion; de lo contrario, el ruido del gradiente es demasiado grande;
    \item usar off-policy correction (como en HIRO) permite corregir el bias del high-level update.
\end{itemize}
\end{warningbox}

\subsection{Hindsight relabeling y reutilización de datos}

El hindsight relabeling toma los subobjetivos no alcanzados del high-level plan y los «reescribe» como un nuevo pseudo-goal a partir de los estados que el nivel bajo sí alcanzó, de modo que el replay buffer aporte más muestras positivas. Al implementarlo, hay que controlar la frecuencia mediante un relabel scheduler para evitar que un rewrite excesivo dañe el original objective; los metadata de todas las relabeled samples deben escribirse en el prompt log para permitir su audit.

\begin{importantbox}{Pasos para implementar el hindsight relabel}
\begin{enumerate}
    \item Recolectar el goal de alto nivel y la trajectory de bajo nivel de un episode;
    \item Elegir un estado como relabel goal;
    \item Recalcular la reward y actualizar la Q de alto nivel;
    \item Escribir los metadata de la relabeled sample (goal hash + timestamp) en la instrumentation;
\end{enumerate}
\end{importantbox}

\subsection{Resumen de la sección}

Las políticas jerárquicas goal-conditioned hacen que el nivel alto fije subobjetivos como un comandante y tratan al nivel bajo como el ejecutor; la dificultad del entrenamiento reside en el drift frecuente del entorno y de los objetivos, lo que exige guardrails de hindsight y off-policy.
\section{Imitación jerárquica y descubrimiento de habilidades}

\subsection{Extracción de habilidades a partir de demostraciones}

El aprendizaje por imitación ya no depende de la reward, sino que infiere la skill segmentation a partir de la demo. Entre los enfoques habituales están sliding window clustering, change point detection y el skill-specific autoencoder.

\begin{importantbox}{Pipeline de descubrimiento de habilidades}
\begin{enumerate}
    \item Etiquetado: dividir la demo en candidate skills;
    \item Entrenar la skill policy: usar el segment como policy context;
    \item Entrenar el scheduler: el nivel alto aprende a predecir la skill sequence;
    \item Deployment: vincular la skill library a un planificador LLN/LLM.
\end{enumerate}
\end{importantbox}

\subsection{Gobernanza de los datos de imitación}

\begin{knowledgebox}{Checklist para gobernar los datos de imitación}
\begin{itemize}
    \item Demo provenance: registrar el id del video o la trayectoria de origen;
    \item Skill annotation: anotar los límites de las subtareas y exportarlos a `labels.json`;
    \item Diversity filter: eliminar los fragmentos de demostración excesivamente repetidos;
    \item Success flag: marcar en cada demostración si se completó, para que la use la low-layer reward.
\end{itemize}
\end{knowledgebox}

\subsection{VQ-VAE y espacio latente discreto}

El VQ-VAE se encarga de comprimir trayectorias variable-length en discrete codebook entries; cada entry corresponde a una habilidad. La política de nivel alto solo necesita predecir el codebook index, y el decoder de nivel bajo puede «decodificar» la secuencia de acciones necesaria.

\begin{figure}[H]
\centering
\includegraphics[width=0.85\textwidth]{slides-images/slide-25.jpg}
\caption{Esquema de la arquitectura de Skill Discovery (slide-25).}
\end{figure}

\begin{warningbox}{Collapse de la variable latente}
Si el codebook size es demasiado pequeño, el VQ-VAE reutiliza una y otra vez la misma habilidad, lo que deja a la política de nivel alto sin diversidad. Las soluciones habituales son añadir una commitment loss o usar EM-style clustering.
\end{warningbox}

\subsection{El lenguaje como habilidad y como comando}

Un LLM puede servir como interface de lenguaje para la policy de nivel alto; el ponente mencionó que SayCan/Inner Monologue, mediante el pipeline language → intent → skills, aumentan la interpretabilidad del ejecutor de nivel bajo.

\begin{knowledgebox}{Puntos clave para gobernar las instrucciones en lenguaje}
\begin{itemize}
    \item El prompt log registra la language input + la expected skill;
    \item Usar un parser para convertir el language en un discrete skill id;
    \item La retroalimentación del nivel bajo (por ejemplo, el success flag) debe escribirse de vuelta en el prompt log, para formar un governance loop.
\end{itemize}
\end{knowledgebox}

\subsection{Resumen de la sección}

El módulo de imitación permite descubrir automáticamente una biblioteca de habilidades a partir de demostraciones; junto con una interfaz de lenguaje/LMM, logra una planificación jerárquica que las personas pueden interpretar.
\section{Planificación de alto nivel guiada por lenguaje y LLM}

\subsection{Los dos papeles del LLM como planificador}

El LLM puede generar una secuencia de subtareas y también puede actuar como feasibility filter. El ponente enfatiza que el plan generado por el LLM debe llevar su provenance (hora de origen, prompt hash) y que cada step debe pasar por un tool-check.

\begin{table}[H]
\centering
\begin{tabular}{p{4cm}p{5cm}}
\toprule
Papel & Puntos prácticos \\
\midrule
Plan generator & Registrar prompt id, temperature y llm version \\
Verifier & Verificar la viabilidad del plan mediante una tool chain (por ejemplo, un python sandbox) \\
Skill mapper & Convertir el plan en skill id o primitive action \\
\bottomrule
\end{tabular}
\caption{Los múltiples papeles del LLM en un flujo jerárquico.}
\end{table}

\begin{importantbox}{LLM Prompt log}
Cada prompt lleva asociada su metadata: `timestamp`, `llm-version`, `temperature`, `tool-call`, `expected skill id`; además, la retroalimentación sobre hallucination/feasibility se devuelve al prompt log.
\end{importantbox}

\subsection{Plan verification pipeline}

\begin{figure}[H]
\centering
\includegraphics[width=0.85\textwidth]{slides-images/slide-33.jpg}
\caption{Flujo de verificación del plan del LLM y de tool-check (slide-33).}
\end{figure}

Cada salida del LLM pasa primero por el verifier (python sandbox + knowledge graph) y después el skill mapper la interpreta como un plan ejecutable. Si se detecta una hallucination, se notifica de inmediato al incident channel y el context se registra en el grade-book.

\begin{warningbox}{Guardrail de plan verification}
\begin{itemize}
    \item Cada plan incluye \texttt{plan\_id}, \texttt{llm\_version} y \texttt{prompt\_hash};
    \item Antes de la ejecución, se verifica la viabilidad con tool-check;
    \item Cuando se detecta una hallucination, se activa un rollback y se corrige el prompt.
\end{itemize}
\end{warningbox}

\subsection{Gobernanza y auditoría}

Al generar la planificación debe haber guardrails de gobernanza, como multi-round verification, tool sandbox y human-in-the-loop review. Si la LLM output contiene una high-risk action (por ejemplo, modify-files), es obligatorio activar un incident drill.

\begin{warningbox}{Riesgos de la planificación con LLM}
La planificación automatizada puede producir hallucination, sobre todo en tareas de varios pasos. Se recomienda establecer, antes del despliegue:
\begin{itemize}
    \item Multi-modal verification: verificar cada paso con una tool call (verify each step);
    \item Stop tokens: evitar que el LLM agregue subtask sin una condición de parada;
    \item Human review: confirmación humana antes de un critical plan.
\end{itemize}
\end{warningbox}

\subsection{Resumen de la sección}

El lenguaje y el LLM dan más explainability al RL jerárquico, pero también trasladan la carga de gobernanza de hallucination/incident a etapas más tempranas: el prompt log de alto nivel y el verification pipeline.
\section{Casos prácticos y gobernanza de ingeniería}

\subsection{Manipulación robótica de horizonte largo}

El ejemplo de la cocina que presenta Chelsea Finn se compone de varias skill y safety guardrail. Dividimos la tarea en: 1) percepción y estimación de estado, 2) ejecución de la skill (por ejemplo, tomar la olla), 3) verification + recovery.

\begin{knowledgebox}{Puntos de monitoreo del Kitchen Robot}
\begin{itemize}
    \item State check: en cada paso se verifica que gripper/contact funcionen con normalidad;
    \item Skill success flag: al terminar, la skill de bajo nivel reporta `success`;
    \item Recovery template: si la skill falla, el nivel alto activa de inmediato un safe fallback;
    \item Human override: antes de una high-risk action (cutting/pouring) se inserta human-in-the-loop.
\end{itemize}
\end{knowledgebox}

\subsection{Matriz de evaluación y monitoreo}

Se introduce un multi-metric scoreboard para dar seguimiento a accuracy, self-consistency, hallucination rate, token budget, etc. En cada emergent jump hay que revisar varias curvas a la vez para no llegar a un juicio erróneo a partir de una single metric.

\begin{table}[H]
\centering
\begin{tabular}{p{3.5cm}p{5cm}p{3cm}}
\toprule
Métrica & Descripción & Herramienta \\
\midrule
Self-consistency & agreement entre generaciones por múltiples rutas & CoT voting harness \\
Hallucination rate & tasa de errores factuales & human eval + KG check \\
Token budget & consumo de recursos & monitoring hook \\
Instruction drift & desviación de versión del Prompt log & versioned dashboard \\
\bottomrule
\end{tabular}
\caption{Matriz de evaluación del sistema jerárquico.}
\end{table}

\begin{warningbox}{Consecuencias de omitir la cadena de evaluación}
\begin{itemize}
    \item Fijarse solo en accuracy puede ocultar hallucination;
    \item El Prompt drift puede invalidar el high-level plan mientras el nivel bajo sigue ejecutándose;
    \item La falta de monitoreo del token budget provoca compute overrun.
\end{itemize}
\end{warningbox}

\subsection{Incident grade-book y lessons}

\begin{table}[H]
\centering
\begin{tabular}{p{4cm}p{6cm}}
\toprule
Campo & Contenido \\
\midrule
Trigger & hallucination/jump/goal failure \\
Root cause & prompt drift / data shift / skill bug \\
Resolution & rollback prompt / retrain options / add guardrail \\
Next action & drill + doc update \\
\bottomrule
\end{tabular}
\caption{Plantilla del incident grade-book.}
\end{table}

\begin{importantbox}{Procedimiento operativo del Grade-book}
\begin{itemize}
    \item Completar en 24 horas la línea de tiempo trigger → analysis → resolution;
    \item Guardar el snapshot de prompt/data/attention en `notes/incident/`;
    \item Actualizar las lessons en `lessons.md` y compartirlas en la weekly retro.
\end{itemize}
\end{importantbox}

\subsection{Resumen de la sección}

La gobernanza de ingeniería requiere encadenar skill execution, LLM planning y la matriz de monitoreo en un ciclo cerrado, y archivar el prompt log + evaluation snapshot en cada incident.
\section{Garantías de despliegue y reproducibilidad}

\subsection{Checklist de preparación para el despliegue}

Antes del despliegue es obligatorio revisar la prompt stability, el token budget, el plan de fallback y los hooks de monitoreo. Se recomienda usar una multi-column table para mostrar el estado actual de cada indicador y escribir el artifact de cada revisión en `deploy/logs/`.

\begin{table}[H]
\centering
\begin{tabular}{p{4cm}p{5cm}p{3cm}}
\toprule
Elemento de revisión & Método de verificación & Estado \\
\midrule
Prompt stability & múltiples prompts + CoT + self-consistency & En monitoreo \\
Token budget & latency + FLOPs shell & Dimensiones alineadas \\
Fallback & offline fallback script + human review & Preconfigurado \\
Monitoring hook & loss/latency/attention drift & hook enabled \\
\bottomrule
\end{tabular}
\caption{Checklist previa al despliegue.}
\end{table}

\begin{importantbox}{Uso de la checklist de preparación}
Antes de cada despliegue se recorre una standardized template: prompt log → evaluation dashboard → fallback script. Si algún elemento no se aprueba, debe registrarse en el incident grade-book y se pospone la puesta en producción.
\end{importantbox}

\subsection{Tooling stack de observabilidad y reproducibilidad}

Al construir el tooling stack se deben incluir la instrumentation, el seguimiento del LLM plan, la trace de las llamadas a herramientas (doc search, sondeo del python executor) y el attention/scoreboard snapshot.

\begin{figure}[H]
\centering
\includegraphics[width=0.85\textwidth]{slides-images/slide-37.jpg}
\caption{Deployment + stack de observabilidad (slide-37).}
\end{figure}

\begin{table}[H]
\centering
\begin{tabular}{p{3.5cm}p{5cm}}
\toprule
Nivel & Herramienta/práctica \\
\midrule
Prompt & Prompt log + versioned dashboard \\
Inference & Monitoring hook + latency/failure metrics \\
Tooling & API call trace + sandbox constraints \\
Observabilidad & Attention/Config snapshot + drift alert \\
\bottomrule
\end{tabular}
\caption{Stack de observabilidad del despliegue.}
\end{table}

\begin{warningbox}{Riesgos de carecer de observabilidad}
\begin{itemize}
    \item Sin attention trace, no es posible confirmar si un emergent jump es real;
    \item sin tool call log, las low-level skills no pueden rastrearse;
    \item la falta de drift alert retrasará el incident trigger.
\end{itemize}
\end{warningbox}

\subsection{Resumen de la sección}

La parte de despliegue y reproducibilidad es la barrera de protección de la práctica: desde la checklist hasta el stack de observabilidad y, después, el incident grade-book, el aim es que los high-level insights puedan llevarse a la práctica.

\section{Acciones y recomendaciones para la puesta en práctica}

\subsection{Lista de acciones a corto plazo}

\begin{table}[H]
\centering
\begin{tabular}{p{3cm}p{6cm}}
\toprule
Ventana de tiempo & Entregables \\
\midrule
Week 1 & Completar data review + skill segmentation; establecer prompt log template \\
Week 2 & Launch evaluation dashboard, registrar multi-metric drift; definir incident grade-book \\
Week 3 & Activar un emergent drill: LLM plan + skill exec y governance review \\
\bottomrule
\end{tabular}
\caption{Acciones recomendadas a corto plazo.}
\end{table}

\begin{importantbox}{Puntos clave para poner en práctica las acciones}
\begin{itemize}
    \item Prompt log + LLN plan deben llevar asociados version/temperature/expected skill;
    \item cada emergent drill debe contar con review notes y dar lugar a lessons learned;
    \item Incident grade-book registra el trigger, la redirect action y el guardrail change posterior.
\end{itemize}
\end{importantbox}

\subsection{Agenda de investigación a largo plazo}

Se recomienda organizar las líneas a largo plazo en forma de quarterly thesis: 1) scale-law + hierarchical optimizer; 2) LLM-driven skill discovery; 3) attention interpretability for subgoal.

\begin{knowledgebox}{Ejemplos de investigación a largo plazo}
\begin{itemize}
    \item Project Alpha: explorar patrones de hierarchical attention;
    \item Project Beta: mapear las curvas de scale law a un optimizer schedule;
    \item Project Gamma: construir un emergent dashboard que apoye las decisiones del governance panel.
\end{itemize}
\end{knowledgebox}

\subsection{Checklist automation}

Se automatiza la action checklist de corto plazo; por ejemplo, con Github Actions se verifica cada semana que el prompt log esté sincronizado y con dashboards se muestra el multi-metric drift. Así se reduce la verificación manual repetitiva y los ingenieros pueden concentrar su esfuerzo en la incident response.

\begin{importantbox}{Recordatorios de la checklist automatizada}
\begin{itemize}
    \item una actualización del prompt activa el workflow \texttt{check\_prompt\_versions};
    \item cuando un drift alert alcanza el threshold, se activa \texttt{incident-drill};
    \item \texttt{summary.md} genera cada semana un resumen de lessons para que lo revise el governance panel.
\end{itemize}
\end{importantbox}

\subsection{Resumen de la sección}

La sección de acciones recomendadas convierte la teoría en un timeline: los drills de corto plazo cierran el ciclo con rapidez, mientras que los proyectos de largo plazo convierten los emergent insights en research output verificable.
\section{Consulta rápida de términos y métricas}

\subsection{Terminología clave}

\begin{figure}[H]
\centering
\includegraphics[width=0.65\textwidth]{slides-images/slide-39.jpg}
\caption{Slide de consulta rápida de términos y métricas (slide-39).}
\end{figure}

\begin{table}[H]
\centering
\begin{tabular}{p{4cm}p{7cm}}
\toprule
Término & Significado \\
\midrule
CoT & Chain-of-Thought: se muestran al modelo los pasos completos del razonamiento \\
Self-consistency & Se muestrean varias rutas y se toma la moda; reduce los errores ocasionales \\
Governance loop & Prompt log + dashboard + incident grade-book \\
KV cache & Caché de key/value; acelera los diálogos de varios turnos en modelos decoder-only \\
\bottomrule
\end{tabular}
\caption{Términos clave que se repiten en la Lecture 15.}
\end{table}

\begin{knowledgebox}{Sugerencia para memorizar los términos}
Escribe los términos clave en `glossary.md` y repásalos en la incident retro; así los ingenieros que se incorporan al equipo se ponen al día con rapidez.
\end{knowledgebox}

\subsection{Matriz de métricas y dashboard}

\begin{table}[H]
\centering
\begin{tabular}{p{4cm}p{5cm}}
\toprule
Métrica & Función \\
\midrule
Self-consistency score & Mide la estabilidad del prompt/plan \\
Hallucination rate & Registra la proporción de eventos de fabrication \\
Token budget & Límite de token/compute por solicitud \\
Attention entropy & Indica si el modelo salta a un modo distinto \\
\bottomrule
\end{tabular}
\caption{Métricas de uso frecuente en la gobernanza y la depuración.}
\end{table}

\begin{warningbox}{Riesgos del retraso de las métricas}
\begin{itemize}
    \item Depender en exceso de la accuracy puede ocultar la hallucination;
    \item si la Hallucination rate sufre label delay, la reacción de governance llega tarde;
    \item sin monitoreo de la attention entropy, es difícil confirmar si un emergent jump es real.
\end{itemize}
\end{warningbox}

\subsection{Resumen de la sección}

Dejar por escrito los términos y las métricas ayuda a que los nuevos integrantes entiendan pronto el vocabulario de la clase, y también permite localizar con rapidez key term + metric en una incident review.
\section{Q\&A y recomendaciones prácticas}

\subsection{Puntos principales de la sesión de preguntas}

En la sesión de Q\&A, Chelsea Finn enfatizó los tres pasos «observe → instrument → govern», en particular la necesidad de identificar el prompt, los datos o el attention artifact concretos después de que ocurre un emergent drift.

\begin{table}[H]
\centering
\begin{tabular}{p{4cm}p{5cm}}
\toprule
Tema de la pregunta & Respuesta central \\
\midrule
Emergent stability & Solo después de que varias métricas coinciden puede declararse madura la emergence \\
Prompt drift & insist on prompt log + rollback plan \\
Tooling stack & Las trazas de attention/plan/metric deben ser realtime + archival \\
\bottomrule
\end{tabular}
\caption{Lo esencial de la sesión de Q\&A.}
\end{table}

\begin{knowledgebox}{Puntos para recordar del Q\&A}
Escribe las keywords del Q\&A en `qa.md' y contrástalas en cada incident review, para asegurarte de no omitir ninguna advertencia de alto riesgo.
\end{knowledgebox}

\subsection{Recomendaciones prácticas}

\begin{figure}[H]
\centering
\includegraphics[width=0.85\textwidth]{slides-images/slide-41.jpg}
\caption{Slide de Q\&A y recomendaciones prácticas (slide-41).}
\end{figure}

\begin{importantbox}{Recomendaciones para la puesta en práctica}
\begin{itemize}
    \item Registrar la prompt version + sampling seed;
    \item cada emergent drill debe producir lessons learned + summary note;
    \item en las reuniones del governance panel, hacer review del incident grade-book + prompt log.
\end{itemize}
\end{importantbox}

\subsection{Resumen de la sección}

El Q\&A convierte los abstract insights de la clase en una checklist; mientras los artifacts que surgen de las preguntas y respuestas se sigan poniendo por escrito, la gobernanza siempre tendrá un fundamento verificable.
\section{Síntesis y ampliación}

\begin{table}[H]
\centering
\begin{tabular}{p{3.5cm}p{8cm}}
\toprule
Dimensión & Revisión central \\
\midrule
Motivación & La estructura jerárquica es fundamental para resolver la asignación de crédito y la exploración \\
Métodos & Options/SMDP + goal-conditioned + imitation (VQ-VAE) ofrecen una ruta completa \\
Gobernanza del lenguaje & LLM prompt log + verification pipeline garantizan que el plan de alto nivel sea confiable \\
Ingeniería | evaluación & scoreboard + incident grade-book permiten rastrear el emergent drift \\
\bottomrule
\end{tabular}
\caption{Repaso de las tres líneas principales de la Lecture 15.}
\end{table}

\begin{table}[H]
\centering
\begin{tabular}{p{3.5cm}p{8cm}}
\toprule
Dimensión & Riesgos y control \\
\midrule
No estacionariedad & hindsight relabel + off-policy correction \\
Hallucination & multi-metric governance + tool check \\
Script drift & versioned prompt log + incident reviews \\
\bottomrule
\end{tabular}
\caption{Perspectiva complementaria sobre riesgos y gobernanza.}
\end{table}

\subsection{Lecturas adicionales}
\begin{enumerate}
    \item Sutton et al., ``Between MDPs and Semi-MDPs: A Framework for Temporal Abstraction in RL,'' 1999
    \item Nachum et al., ``Data-Efficient Hierarchical Reinforcement Learning (HIRO),'' NeurIPS 2018
    \item Pertsch et al., ``Accelerating RL with Learned Skill Priors (SPiRL),'' CoRL 2020
    \item Ahn et al., ``Do As I Can, Not As I Say: Grounding Language in Robotic Affordances (SayCan),'' 2022
    \item Ajay et al., ``Compositional Foundation Models for Hierarchical Planning,'' NeurIPS 2023
\end{enumerate}

\end{document}
